% Fragment -- inclus de main.tex prin \input
\clearpage
\setpartlabel{themeDeploy}{\faRocket}{Faza 4: Deploy}
\setcounter{section}{0}
\phantomsection
\addcontentsline{toc}{part}{\protect\textbf{Faza 4 \textendash{} Deploy}}
\handout{Faza 4 \textendash{} Deploy}{Versionare cu git. Publicare prin GitHub Pages.}

Capitolul descrie procesul de publicare a unui site web static prin combinatia git \textendash{} GitHub \textendash{} GitHub Pages. Acoperirea include comenzile esentiale ale fluxului de lucru, configurarea unui repository remote, activarea serviciului de hosting integrat, si convenii de scriere a mesajelor de commit.

\bigskip

\begin{outcomes}
Dupa capitolul asta vei \textcommabelow{s}ti:
\begin{itemize}[itemsep=1pt]
  \item Executa cele 5 comenzi git pentru ini\textcommabelow{t}ializarea unui repository \textcommabelow{s}i primul push.
  \item Configureaza GitHub Pages pe un repository pentru hosting static automat.
  \item Folose\textcommabelow{s}te \code{git status}, \code{git diff}, \code{git log} pentru inspectarea starii.
  \item Scrie mesaje de commit la imperativ, descriptive, sub 72 caractere.
  \item Excludre fi\textcommabelow{s}iere sensibile prin \code{.gitignore}.
\end{itemize}
\end{outcomes}

\begin{whymatters}
Site-ul construit local pe ma\textcommabelow{s}ina dezvoltatorului nu este accesibil nimanui altcuiva. Deploy-ul transforma cod local in URL public. In plus, git este standardul industrial pentru versionare \textendash{} cuno\textcommabelow{s}tin\textcommabelow{t}ele dobandite aici se aplica identic in orice job de dezvoltare software. Mesajele de commit bune sunt diferen\textcommabelow{t}a intre un istoric care ajuta la debug \textcommabelow{s}i unul inutil.
\end{whymatters}

\begin{nota}[Pre-rechizite]
Capitolul presupune un site func\textcommabelow{t}ional in localhost (Faza 1\textendash{}3) \textcommabelow{s}i un cont GitHub creat in prealabil.
\end{nota}

% =========================================================
\section{De ce git?}

Git este sistemul de control al versiunilor (\emph{Version Control System}) folosit de aproape toata industria software. Beneficiile fata de o solutie de tip ,,upload manual prin FTP'':

\begin{itemize}
  \item \textbf{Istoric.} Fiecare commit reprezinta un \emph{snapshot} al codului, identificat printr-un hash unic. Reverenta la o versiune anterioara este o operatie atomica.
  \item \textbf{Backup distribuit.} Repository-ul exista local, pe GitHub, si pe orice copie clonata. Pierderea unei singure copii nu afecteaza istoricul.
  \item \textbf{Colaborare.} Mai multi autori pot lucra simultan pe ramuri (\emph{branches}) separate, sincronizate prin merge.
  \item \textbf{Standard industrial.} Cunoasterea git este o cerinta de baza in posturile de dezvoltare software.
\end{itemize}

% =========================================================
\section{Modelul mental: trei zone}

Git mentine trei zone in care fisierele pot exista la un moment dat:

\begin{lstlisting}[style=base, basicstyle=\ttfamily\small, numbers=none]
  Working dir         Staging area         Repository
  (fisiere editate)   (pregatite pt.       (snapshot-uri
                       commit)              salvate permanent)

       |                   |                    |
       |  git add .        |  git commit -m     |
       |------------------>|------------------->|
       |                   |                    |
       |   git restore     |    git reset       |
       |<------------------|<-------------------|
       |                   |                    |
                                                |
                                                |  git push
                                                |--------->  GitHub
                                                |             (remote)
                                                |<---------
                                                |  git pull
\end{lstlisting}

\begin{itemize}
  \item \textbf{Working directory} \textendash{} fisierele asa cum apar pe disc.
  \item \textbf{Staging area} (sau \emph{index}) \textendash{} multimea fisierelor pregatite pentru urmatorul commit.
  \item \textbf{Repository} \textendash{} istoricul snapshot-urilor confirmate.
\end{itemize}

Comenzile \code{add}, \code{commit}, \code{push} muta starea intre zone in directia stanga $\to$ dreapta.

% =========================================================
\section{Cele cinci comenzi de baza}

Fluxul minim pentru publicarea unui director nou pe GitHub:

\begin{lstlisting}[style=bash]
git init -b main                               # 1
git add .                                      # 2
git commit -m "Initial commit"                 # 3
git remote add origin URL_DE_PE_GITHUB         # 4
git push -u origin main                        # 5
\end{lstlisting}

\begin{itemize}
  \item \code{git init -b main} \textendash{} initializeaza un repository local cu ramura implicita \code{main}.
  \item \code{git add .} \textendash{} include toate fisierele modificate in zona de pregatire (\emph{staging area}).
  \item \code{git commit -m "..."} \textendash{} creeaza un snapshot al starii pregatite, cu un mesaj descriptiv.
  \item \code{git remote add origin ...} \textendash{} asociaza repository-ul local cu un remote pe GitHub.
  \item \code{git push -u origin main} \textendash{} trimite commit-urile locale catre remote, stabilind ramura \code{main} ca \emph{upstream}.
\end{itemize}

% =========================================================
\section{Comenzi de inspectie (read-only)}

Inainte de \code{commit}, doua comenzi confirma ce urmeaza sa fie salvat:

\begin{lstlisting}[style=bash, caption={Verificare inainte de commit}]
git status                  # Lista fisierelor modificate / pregatite
git diff                    # Diferenta linie cu linie (working dir)
git diff --staged           # Diferenta in staging area
git log --oneline -10       # Ultimele 10 commit-uri, format compact
\end{lstlisting}

Aceste comenzi nu modifica nimic; pot fi rulate oricand fara risc.

% =========================================================
\section{Configurarea GitHub Pages}

Procedura standard pentru activarea hosting-ului:

\begin{enumerate}
  \item Creare a unui repository public la \url{https://github.com/new}, fara fisiere implicite (README, gitignore).
  \item Copierea URL-ului HTTPS afisat de GitHub si executarea comenzii \code{git remote add origin <URL>} urmata de \code{git push -u origin main}.
  \item Navigare la \emph{Settings} $\to$ \emph{Pages} (in panoul lateral al repository-ului).
  \item Selectarea ramurii \code{main} si a folder-ului \code{/ (root)} in sectiunea \emph{Branch}, urmata de \emph{Save}.
  \item Asteptarea de 30 \textendash{} 60 secunde, dupa care URL-ul public devine accesibil in caseta verde din partea de sus.
\end{enumerate}

\textbf{Anatomia URL-ului public:} \code{https://USERNAME.github.io/REPO}

% =========================================================
\section{Workflow zilnic dupa primul push}

Actualizarea continutului dupa deploy-ul initial necesita doar trei comenzi:

\begin{lstlisting}[style=bash]
git status                              # ce s-a modificat?
git add .                               # marcare pentru commit
git commit -m "Adauga sectiunea proiecte"
git push                                # trimitere pe GitHub
\end{lstlisting}

GitHub Pages reconstruieste si republica pagina automat, in aproximativ 60 de secunde.

% =========================================================
\section{Convenii pentru mesajele de commit}

Mesajul de commit ramane in istoricul repository-ului permanent. Calitatea acestuia este vizibila la \code{git log} si afecteaza utilitatea istoricului ca documentatie a evolutiei proiectului.

\begin{dano}[Asa nu]
\begin{lstlisting}[style=bash, numbers=none]
git commit -m "update"
git commit -m "fix"
git commit -m "asdf"
git commit -m "Modificari"
git commit -m "."
\end{lstlisting}
\textbf{Probleme.} Niciun context. La revizuire, autorul nu mai stie ce contine fiecare commit. Inutilizabil pentru \code{git revert}, \code{git bisect}, code review.
\end{dano}

\begin{dado}[Asa da]
\begin{lstlisting}[style=bash, numbers=none]
git commit -m "Adauga sectiunea hero cu CTA"
git commit -m "Fix overflow horizontal pe mobil < 360px"
git commit -m "Implementeaza dark mode prin light-dark()"
git commit -m "Reduce dimensiunea hero.jpg de la 2.4MB la 180KB"
\end{lstlisting}
\textbf{Reguli simple.} (1) Verb la imperativ in prima pozitie (\emph{Adauga}, \emph{Fix}, \emph{Implementeaza}, \emph{Reduce}). (2) Maxim 50\textendash{}72 caractere. (3) Descrie \emph{ce} si \emph{de ce}, nu \emph{cum} (codul arata cum).
\end{dado}

% =========================================================
\section{Fisierul .gitignore}

Anumite fisiere si directoare nu trebuie sa ajunga in repository: dependinte instalate (\code{node\_modules/}), fisiere temporare ale editoarelor, secretele (\code{.env}). Lista lor se declara in fisierul \code{.gitignore} la radacina proiectului:

\begin{lstlisting}[style=bash, caption={Exemplu minimal pentru proiect web static}]
# Dependinte
node_modules/

# Editoare
.vscode/
.idea/
*.swp

# Sistem
.DS_Store
Thumbs.db

# Secrete
.env
.env.local

# Build output
dist/
build/
\end{lstlisting}

Fisierul se commit-eaza ca orice alt fisier. Toate caile listate sunt ignorate la \code{git add}.

\begin{atentie}
\textbf{Nu commite niciodata fisiere \code{.env} cu chei API, parole, token-uri.} Odata commit-ate, raman in istoric chiar dupa stergere; pot fi recuperate de oricine cu acces la repository. Daca s-a intamplat: regenerati cheia compromisa imediat la furnizorul respectiv.
\end{atentie}

% =========================================================
\section{Probleme frecvente}

\begin{atentie}
\textbf{Autentificare prin parola.} GitHub a renuntat la autentificarea prin parola pentru operatiunile git (din 2021). Solutia recomandata este generarea unui \emph{Personal Access Token} (PAT) la \url{https://github.com/settings/tokens} cu scope-ul \code{repo}, folosit ulterior in locul parolei. Browserul memoreaza tokenul pentru sesiunile viitoare.
\end{atentie}

\begin{atentie}
\textbf{Eroare 404 pe pagina deployata.} Cauze posibile: (1) intervalul de propagare nu a fost respectat (60 secunde dupa \emph{Save}); (2) fisierul de pornire nu se numeste \code{index.html} (este sensibil la majuscule pe serverul GitHub); (3) cache-ul browserului afiseaza versiunea precedenta \textendash{} se rezolva cu \code{Ctrl+F5}.
\end{atentie}

% =========================================================
\section*{Verificare cuno\textcommabelow{s}tin\textcommabelow{t}e}

\begin{selfcheck}
\begin{enumerate}[itemsep=2pt]
  \item Care sunt cele 5 comenzi git pentru ini\textcommabelow{t}ializare \textcommabelow{s}i primul push?
  \item De ce GitHub a renun\textcommabelow{t}at la autentificarea prin parola? Care este alternativa?
  \item \emph{Spot the issue:} \code{git commit -m "update"}.
  \item Diferen\textcommabelow{t}a intre \code{git status} \textcommabelow{s}i \code{git diff}?
  \item Pentru un repository numit \code{portfolio} al utilizatorului \code{maria}, care este URL-ul GitHub Pages rezultat?
\end{enumerate}
\end{selfcheck}

% =========================================================
\section*{Recapitulare}

\begin{recap}[Faza 4 -- Deploy]
\begin{itemize}[itemsep=1pt]
  \item Modelul mental: working dir $\to$ staging $\to$ repository $\to$ remote.
  \item Cinci comenzi de baza: \code{git init -b main}, \code{git add .}, \code{git commit -m}, \code{git remote add origin}, \code{git push -u origin main}.
  \item Comenzi de inspectie: \code{git status}, \code{git diff}, \code{git log --oneline -10}.
  \item Activare Pages: \emph{Settings} $\to$ \emph{Pages} $\to$ branch \code{main}, folder \code{/ (root)} $\to$ \emph{Save}.
  \item Mesaje de commit: verb la imperativ, $\leq$72 caractere, descrie \emph{ce} si \emph{de ce}.
  \item Workflow zilnic: 3 comenzi (\code{add}, \code{commit}, \code{push}); redeploy automat in $\sim$60s.
  \item Autentificare prin Personal Access Token (PAT), nu prin parola.
  \item \code{.gitignore} exclude \code{node\_modules/}, \code{.env}, fisiere de editor, build artifacts.
\end{itemize}
\end{recap}

\partdivider

% =========================================================
\begin{joc}
\textbf{Platforma:} GitHub (\url{https://github.com}) \textendash{} aplicatia practica este deploy-ul efectiv al site-ului construit in fazele 1\textendash{}3.

\medskip

\textbf{Activitatea.} Procedura completa, executata sincron de toti participantii sub indrumarea trainer-ului:
\begin{enumerate}[itemsep=1pt]
  \item Inchiderea fisierelor in editor; navigare in \emph{Git Bash} la directorul proiectului.
  \item Executarea celor cinci comenzi git din sectiunea ,,Cele cinci comenzi de baza''.
  \item Crearea repository-ului public pe GitHub si conectarea remote-ului.
  \item Activarea GitHub Pages in setari.
  \item Verificarea URL-ului public si publicarea acestuia in canalul Discord al grupului.
\end{enumerate}

\textbf{Durata estimata:} 10 minute.
\end{joc}

% =========================================================
\section*{\faGraduationCap\ \ Resurse pentru aprofundare}

\begin{itemize}[itemsep=2pt]
  \item GitHub Inc. \emph{GitHub Pages documentation}. \url{https://pages.github.com}
  \item Cottle P. \emph{Learn Git Branching}. \url{https://learngitbranching.js.org}
  \item Chacon S., Straub B. \emph{Pro Git} (editia a 2-a). \url{https://git-scm.com/book/en/v2}
  \item GitHub Inc. \emph{Configuring a custom domain for your GitHub Pages site}. \\\url{https://docs.github.com/pages/configuring-a-custom-domain-for-your-github-pages-site}
  \item Beams C. \emph{How to Write a Git Commit Message}. \url{https://cbea.ms/git-commit/}
  \item Alternative: \url{https://www.netlify.com}, \url{https://vercel.com} (servicii similare cu deploy automat din git).
\end{itemize}

\bigskip

\bridge{Capitolul 5 (SEO) face site-ul descoperibil. URL public exista; acum trebuie ca motoarele de cautare \textcommabelow{s}i re\textcommabelow{t}elele sociale sa-l afi\textcommabelow{s}eze corect.}
